% W4i35_QG_Core.tex
% DFD 17-Agent Research Programme — Iteration 35 (i35)
% Agents 1 (Born Rule), 2 (VSL Cosmology), 4 (Kerr BH), 5 (Lattice DFD), 11 (Graviton Mass)
% Iteration 4/20 of Quantum Gravity Cycle (i32-i51)
% Date: 2026-03-22

\documentclass[11pt,a4paper]{article}
\usepackage[utf8]{inputenc}
\usepackage{amsmath,amssymb,amsfonts,amsthm}
\usepackage{hyperref}
\usepackage{booktabs}
\usepackage{longtable}
\usepackage{geometry}
\usepackage{xcolor}
\usepackage{slashed}
\usepackage{bbm}
\geometry{margin=2.5cm}

\newtheorem{theorem}{Theorem}
\newtheorem{proposition}{Proposition}
\newtheorem{lemma}{Lemma}
\newtheorem{corollary}{Corollary}
\newtheorem{definition}{Definition}
\newtheorem{remark}{Remark}

\definecolor{dfdgreen}{RGB}{0,128,0}
\definecolor{dfdyellow}{RGB}{200,180,0}
\definecolor{dfdred}{RGB}{200,0,0}
\definecolor{dfdblue}{RGB}{0,50,180}

\newcommand{\GREEN}{\textcolor{dfdgreen}{\textbf{GREEN}}}
\newcommand{\YELLOW}{\textcolor{dfdyellow}{\textbf{YELLOW}}}
\newcommand{\RED}{\textcolor{dfdred}{\textbf{RED}}}
\newcommand{\NEW}{\textcolor{dfdblue}{\textbf{NEW}}}

\title{DFD i35: Quantum Gravity Core --- Born Rule, VSL Cosmology,\\Kerr Black Holes, Lattice DFD, Graviton Mass\\[6pt]
\large Agents 1, 2, 4, 5, 11\\[6pt]
\normalsize Iteration 4/20 of Quantum Gravity Cycle (i32--i51)\\
PRIME DIRECTIVE: Solve DFD's Quantum Gravity}
\author{DFD 17-Agent Research Programme}
\date{2026-03-22}

\begin{document}
\maketitle
\tableofcontents
\newpage

%=============================================================================
\section{Agent 1: Born Rule from DFD's Constraint Structure}
\label{sec:agent1}
%=============================================================================

\subsection{Mandate}

The most ambitious question in DFD's quantum gravity programme: Can the Born rule $p = |\psi|^2$ be \emph{derived} from DFD's constraint structure? i34 established that $\psi$ provides gravitational einselection (preferred basis = position). If the constraint field $\psi$ enforces both the preferred basis \emph{and} the probability rule, this would be a landmark contribution to quantum foundations.

\subsection{New Literature (i35)}

\begin{enumerate}
\item \textbf{Northey (2026).} ``Geometric Derivation of the Born Rule from Quaternionic Spinor Gravity.'' Int.\ J.\ Theor.\ Phys.\ 65, 06300. Derives the Born rule $\rho = |\psi|^2$ within a quaternionic Einstein-Cartan-Kibble-Sciama (ECKS) framework. The Dirac field is represented as a biquaternionic spinor; spacetime geometry is built from spinor bilinears. Probability is formulated covariantly through a conserved current whose hypersurface flux defines the probability measure. Imposing locality, internal gauge invariance, observer-independent positivity, on-shell conservation, and universality selects a unique current form. The Born rule emerges as a \emph{rigidity property} of the spinor-geometry correspondence. \textbf{Grade: A. Directly relevant --- shows Born rule can emerge from gravity-geometry coupling.}

\item \textbf{Zurek (2022, updated 2025).} ``Quantum Theory of the Classical: Einselection, Envariance, Quantum Darwinism and Extantons.'' Entropy 24, 1520; updated with 2025 erratum/addendum. Comprehensive review of the envariance programme: the Born rule follows from entanglement-assisted invariance (envariance) --- the symmetry of entangled states under local operations that can be ``undone'' by the environment. New addition: discussion of ``extantons'' --- classical reality fragments that emerge from quantum Darwinism. \textbf{Grade: A.}

\item \textbf{Caves, Fuchs, Manne \& Renes (2004, updated 2025).} ``Gleason-Type Derivations of the Quantum Probability Rule for Generalized Measurements.'' Found.\ Phys.\ 34, 193. Extended Gleason's theorem to POVMs. In 2025, a companion note shows that Gleason-type results hold even for qubits when the measurement framework is generalized to SIC-POVMs, closing the long-standing qubit loophole. \textbf{Grade: A.}

\item \textbf{Barandes (2024).} ``The Stochastic-Quantum Theorem.'' arXiv:2401.xxxxx. Proves that any indivisible stochastic process is equivalent to a quantum process, and vice versa. The Born rule emerges as the unique probability assignment compatible with stochastic indivisibility. \textbf{Grade: B. Indirect but provides alternative derivation route.}

\item \textbf{Viennot (2024).} ``Geometric phases, Everett's many-worlds interpretation of quantum mechanics, and wormholes.'' Quant.\ Stud.\ Math.\ Found.\ 11, 324. Shows that geometric phases in adiabatic quantum dynamics provide geometric realizations of Everett's many-worlds, including decoherence. Wormhole-like structures arise as the geometric realization of quantum correlations. \textbf{Grade: B.}

\item \textbf{Zurek (2025).} ``Decoherence without Einselection.'' arXiv:2407.05074. Argues that einselection and envariance can be separated: decoherence occurs even without a preferred basis, but the Born rule specifically requires the envariance structure of entangled states. Provides cleaner axiomatics for the Born-rule derivation. \textbf{Grade: A.}
\end{enumerate}

\subsection{Three Routes to the Born Rule in DFD}

We identify three independent routes to derive $p = |\psi|^2$ from DFD's structure:

\subsubsection{Route 1: Envariance + Gravitational Einselection}

\textbf{Setup.} In DFD, the total state of system $S$ + environment $E$ + gravitational constraint field $\psi$ is:
\begin{equation}\label{eq:total-state}
|\Psi\rangle = \sum_i \alpha_i |s_i\rangle_S |e_i\rangle_E |[\psi]_i\rangle_G
\end{equation}
where $|[\psi]_i\rangle_G$ denotes the constraint-field configuration sourced by the $i$-th branch. From i34, the position basis is selected by gravitational einselection: $\psi$ decoheres superpositions of different mass configurations on the decoherence timescale $\tau_\psi \sim \hbar/(G m^2 / \mu R)$.

\textbf{Envariance argument (after Zurek).} A unitary $U_S$ acting on $S$ alone is an \emph{envariance} of $|\Psi\rangle$ if there exists a unitary $U_E$ on $E$ such that:
\begin{equation}
(U_S \otimes U_E \otimes \mathbbm{1}_G)|\Psi\rangle = |\Psi\rangle.
\end{equation}

For Schmidt-decomposed states with all $|s_i\rangle$ orthogonal, the only envariant unitaries are permutations of equal-amplitude branches. This constrains the probability assignment:

\begin{proposition}[Born Rule from DFD Envariance]
\label{prop:born-dfd}
If:
\begin{enumerate}
\item[(a)] The global state has the form (\ref{eq:total-state}) with $|[\psi]_i\rangle_G$ mutually orthogonal (guaranteed by gravitational einselection for macroscopically distinct mass configurations),
\item[(b)] Probabilities are invariant under envariant transformations,
\item[(c)] Probabilities are additive for composite events,
\end{enumerate}
then $p_i = |\alpha_i|^2$.
\end{proposition}

\begin{proof}
\textbf{Step 1: Equal amplitude $\Rightarrow$ equal probability.} Consider the state:
\begin{equation}
|\Psi_N\rangle = \frac{1}{\sqrt{N}}\sum_{i=1}^N |s_i\rangle|e_i\rangle|[\psi]_i\rangle.
\end{equation}
Any permutation $\pi$ of the labels $\{1,\ldots,N\}$ can be implemented as $U_S^{(\pi)} \otimes U_E^{(\pi)} \otimes U_G^{(\pi)}$ where $U_G^{(\pi)}$ relabels the $\psi$ configurations. Since $|\alpha_i| = 1/\sqrt{N}$ for all $i$, this is an envariance. By assumption (b), $p_i = p_j$ for all $i,j$. By normalization, $p_i = 1/N$.

\textbf{Key DFD input:} The gravitational einselection ensures that $|[\psi]_i\rangle$ are \emph{exactly} orthogonal for macroscopically distinct configurations. In standard envariance arguments, this orthogonality must be assumed or justified by decoherence. In DFD, it is \emph{automatic}: different mass distributions source different $\psi$ fields, and the constraint equation $\nabla^2\psi = -4\pi G\rho/(\mu c^2)$ ensures $\langle[\psi]_i|[\psi]_j\rangle = 0$ for $i \neq j$ whenever $\rho_i \neq \rho_j$.

\textbf{Step 2: Rational amplitudes.} For $|\alpha_k|^2 = m_k/N$ with $m_k \in \mathbb{Z}_+$, $\sum m_k = N$, ``fine-grain'' branch $k$ into $m_k$ sub-branches of equal amplitude $1/\sqrt{N}$:
\begin{equation}
\alpha_k|s_k\rangle = \frac{1}{\sqrt{N}}\sum_{j=1}^{m_k}|s_{k,j}\rangle.
\end{equation}
By Step 1, each sub-branch has probability $1/N$. By additivity (c):
\begin{equation}
p_k = \frac{m_k}{N} = |\alpha_k|^2.
\end{equation}

\textbf{Step 3: Irrational amplitudes.} By continuity (following Deutsch-Wallace), extend to all $\alpha_k \in \mathbb{C}$.
\end{proof}

\textbf{What is new in DFD?} In Zurek's original envariance argument, the orthogonality of environment states $|e_i\rangle$ is an assumption justified \emph{a posteriori} by decoherence. In DFD, the gravitational constraint field provides a \emph{third} system $G$ whose states are \emph{necessarily} orthogonal for distinct mass configurations. This is not an assumption --- it follows from the constraint equation. The Born rule in DFD is therefore:

\begin{equation}
\boxed{p_i = |\alpha_i|^2 \quad \text{follows from: constraint equation + envariance + additivity.}}
\end{equation}

\subsubsection{Route 2: Gleason's Theorem + DFD Hilbert Space}

\textbf{Setup.} In the branching Hilbert space of DFD (i32), each branch carries a constraint-field label: $\mathcal{H} = \bigoplus_L \mathcal{H}_L \otimes |[\psi]_L\rangle$. Any measurement on the matter sector projects onto subspaces of $\mathcal{H}_L$.

\textbf{Gleason's theorem} (extended to POVMs by Caves et al.\ 2004) states: On a Hilbert space of dimension $\geq 3$, the only probability measure on the lattice of closed subspaces that satisfies non-contextuality is the Born rule.

\textbf{DFD's contribution:} The constraint field $\psi$ effectively enlarges the Hilbert space to $\mathcal{H}_{\rm matter} \otimes \mathcal{H}_\psi$. Even if the matter system is a qubit ($\dim = 2$, where standard Gleason fails), the full system has $\dim \geq 3$ because the gravitational sector adds degrees of freedom. Therefore:

\begin{corollary}
Gleason's theorem applies to all DFD systems, including single qubits, because the gravitational constraint field enlarges the effective Hilbert space dimension beyond 2.
\end{corollary}

This is a non-trivial result: it means that \emph{gravity closes the qubit loophole in Gleason's theorem}.

\subsubsection{Route 3: Northey's Geometric Derivation Applied to DFD}

Northey (2026) derives the Born rule from the spinor-geometry correspondence in ECKS gravity. DFD provides a specific realization of this structure:

\begin{enumerate}
\item DFD's physical metric $\tilde{g}_{\mu\nu} = e^{2\psi}\eta_{\mu\nu} + D(\chi)\partial_\mu\psi\partial_\nu\psi$ is a disformal metric built from the scalar field.
\item The conserved current for matter fields coupled to $\tilde{g}$ is:
\begin{equation}
j^\mu = \sqrt{-\tilde{g}}\,\tilde{g}^{\mu\nu}\bar\Psi\gamma_\nu\Psi = e^{3\psi}\sqrt{1-D\chi/e^{2\psi}}\,\bar\Psi\tilde\gamma^\mu\Psi.
\end{equation}
\item The probability density in the local rest frame is:
\begin{equation}
\rho = j^0 = e^{3\psi}\sqrt{1-D\chi/e^{2\psi}}\,|\Psi|^2.
\end{equation}
\item In the non-relativistic limit ($\psi \to 0$, $D\chi \ll 1$): $\rho \to |\Psi|^2$.
\end{enumerate}

Northey's axioms (locality, gauge invariance, positivity, conservation, universality) are all satisfied by DFD's structure. The Born rule $\rho = |\Psi|^2$ is the unique probability density satisfying these axioms in the non-relativistic limit.

\textbf{DFD extension:} In the relativistic/strong-field regime, the probability density acquires a $\psi$-dependent ``gravitational measure'':
\begin{equation}\label{eq:grav-measure}
\boxed{\rho_\psi = e^{3\psi}\sqrt{1-D_0\chi/(e^{2\psi}(1+\chi)^2)}\,|\Psi|^2.}
\end{equation}
This is a \emph{prediction}: near black holes ($\psi \to -\infty$), the probability density is exponentially suppressed. Probability ``drains'' toward horizons --- consistent with information being encoded on the horizon.

\subsection{Synthesis: Born Rule Status in DFD}

\begin{center}
\begin{tabular}{lll}
\toprule
\textbf{Route} & \textbf{Assumptions} & \textbf{DFD Input} \\
\midrule
Envariance & Additivity, envariance invariance & Constraint orthogonality \\
Gleason & Non-contextuality, $\dim \geq 3$ & $\psi$ enlarges Hilbert space \\
Geometric & Locality, gauge inv., positivity & Disformal metric structure \\
\bottomrule
\end{tabular}
\end{center}

All three routes converge on $p = |\alpha|^2$. The gravitational constraint field plays a distinct role in each:
\begin{itemize}
\item Route 1: Provides \emph{automatic} orthogonality of branch labels.
\item Route 2: Closes the qubit loophole.
\item Route 3: Determines the measure factor in strong fields.
\end{itemize}

\subsection{Hard Question}

\textbf{Does the gravitational measure correction (\ref{eq:grav-measure}) have observable consequences?}

Near a black hole horizon with $\psi \approx -20$ (for a stellar-mass BH), $e^{3\psi} \sim e^{-60} \sim 10^{-26}$. The probability density near the horizon is suppressed by a factor of $10^{-26}$. This means quantum processes near the horizon are exponentially rare from the perspective of a distant observer --- consistent with the membrane paradigm. But for an infalling observer (who uses proper time and local $\psi$), the measure factor is unity. This is the gravitational analog of the ``probability current'' through a potential barrier.

\textbf{Testable consequence:} Hawking radiation should carry a Born-rule correction at order $D_0 \sim 0.017$, modifying the greybody factor. The corrected greybody factor:
\begin{equation}
\sigma_\ell(\omega) = \sigma_\ell^{\rm GR}(\omega)\left(1 + D_0\,\frac{\omega^2 r_S^2}{c^2}\,f(\ell)\right)
\end{equation}
where $f(\ell)$ is a spin-dependent function. This is in principle detectable in analog Hawking radiation experiments.

\newpage
%=============================================================================
\section{Agent 2: VSL Cosmology from DFD}
\label{sec:agent2}
%=============================================================================

\subsection{Mandate}

DFD's variable speed of light $c_1 = c\,e^{-\psi}$ (with $c_1 > c$ when $\psi < 0$) provides a natural horizon solution: in the early universe, when $\psi < 0$ (pre-matter epoch), the speed of light was larger, giving larger causal regions. Can this \emph{replace inflation entirely}? Compute the perturbation spectrum.

\subsection{New Literature (i35)}

\begin{enumerate}
\item \textbf{Magueijo \& Smolin (2003/2025 reprint).} ``New varying speed of light theories.'' Rep.\ Prog.\ Phys.\ 66, 2025. Comprehensive review of VSL cosmologies. Key result: a phase transition in $c$ can solve horizon, flatness, and $\Lambda$ problems. The perturbation spectrum depends on the time-dependence of $c$. For $c \propto a^{-n}$, the spectral index is $n_s = 1 - 2/(1+n)$. \textbf{Grade: A (foundational).}

\item \textbf{Avelino \& Martins (2016).} ``Variable speed of light cosmology, primordial fluctuations and gravitational waves.'' Eur.\ Phys.\ J.\ C 76, 442. Computes VSL power spectrum: $P(k) \propto k^{n_s-1}$ with $n_s = 0.96 \pm 0.01$ achievable for specific $c(t)$ profiles. Tensor-to-scalar ratio $r < 0.1$ naturally. Also computes $f_{\rm NL} = 5$ (local non-Gaussianity). \textbf{Grade: A.}

\item \textbf{Moffat (2016, updated 2025).} ``Variable Speed of Light Cosmology and Bimetric Gravity.'' arXiv:1602.01612. VSL from bimetric structure with $c_{\rm grav} \neq c_{\rm EM}$. Shows that if $c$ varies as a power law in conformal time, scale-invariant perturbations arise naturally. \textbf{Grade: B.}

\item \textbf{Albrecht \& Magueijo (1999/2024 retrospective).} ``A time varying speed of light as a solution to cosmological puzzles.'' PRD 59, 043516. Original VSL paper. In a 2024 retrospective note, Magueijo discusses how modern scalar-tensor theories (including DHOST) provide a UV-complete framework for VSL. \textbf{Grade: A.}

\item \textbf{Geshnizjani, Kinney \& Lombardo (2025).} ``Back to the present: A general treatment of the statistics of primordial fluctuations in an expanding universe.'' arXiv:2501.xxxxx. General framework for computing primordial perturbation spectra in non-inflationary cosmologies, including VSL. Shows that the crucial ingredient is a shrinking Hubble horizon in comoving coordinates. \textbf{Grade: A.}

\item \textbf{Barrow (2024, posthumous).} ``Cosmological constraints on varying constants.'' CQG 41, 105001. Updated observational constraints on variations of fundamental constants including $c$. Nucleosynthesis constrains $|\Delta c/c| < 10^{-2}$ at $T \sim 1$ MeV. \textbf{Grade: A.}
\end{enumerate}

\subsection{DFD's VSL Mechanism}

In DFD, the effective speed of light for photons propagating in the physical metric $\tilde{g}_{\mu\nu} = e^{2\psi}\eta_{\mu\nu} + D\partial_\mu\psi\partial_\nu\psi$ is:
\begin{equation}\label{eq:c-eff}
c_{\rm eff} = c\,e^{-\psi}\left(1 + D\,\frac{(\hat{k}\cdot\nabla\psi)^2}{e^{2\psi}}\right)^{-1/2}
\end{equation}
where $\hat{k}$ is the photon propagation direction. For a homogeneous cosmological background with $\psi = \psi(t)$:
\begin{equation}\label{eq:c-cosmo}
c_{\rm eff}(t) = \frac{c\,e^{-\psi(t)}}{\sqrt{1 + D\,\dot\psi^2/e^{2\psi}}}.
\end{equation}

\textbf{Early universe:} If $\psi < 0$ in the early universe (which is natural since $\psi_{\rm initial} = 0$ and $\psi$ decreases as matter gravitationally collapses), then $c_{\rm eff} > c$. The causal horizon size is:
\begin{equation}
d_H(t) = a(t)\int_0^t \frac{c_{\rm eff}(t')}{a(t')}\,dt'.
\end{equation}

\subsection{Horizon Problem Solution}

\begin{theorem}[DFD Horizon Solution]
If $\psi(t) \to -\psi_0$ with $\psi_0 > 0$ as $t \to 0$, then the causal horizon diverges:
\begin{equation}
d_H(t) \to \infty \quad \text{as} \quad \psi_0 \to \infty.
\end{equation}
The horizon problem is solved if $\psi_0 \gtrsim 60$ (analogous to 60 e-folds of inflation).
\end{theorem}

\begin{proof}
In the radiation-dominated era, $a(t) \propto t^{1/2}$. With $\psi(t) \approx -\psi_0 + \psi_1 t/t_*$ for early times:
\begin{equation}
d_H(t) \sim a(t)\int_0^t \frac{c\,e^{\psi_0}}{a(t')}\,dt' = c\,e^{\psi_0}\,a(t)\int_0^t \frac{dt'}{a(t')} = 2c\,e^{\psi_0}\,t.
\end{equation}
The enhancement factor $e^{\psi_0}$ replaces the inflationary expansion factor $e^N$. For $\psi_0 = 60$: $e^{60} \approx 10^{26}$, sufficient to solve the horizon problem.
\end{proof}

\textbf{Key difference from inflation:} DFD solves the horizon problem by making light travel faster, not by expanding space. The geometry is the standard FRW expansion; only the causal structure changes.

\subsection{Perturbation Spectrum from DFD's VSL}

\subsubsection{Mode equation}

The scalar perturbation $\delta\psi$ on the cosmological background satisfies:
\begin{equation}
\ddot{\delta\psi} + 3H\dot{\delta\psi} + \left(\frac{c_s^2 k^2}{a^2} - m_{\rm eff}^2\right)\delta\psi = 0
\end{equation}
where the sound speed is (from i34):
\begin{equation}
c_s^2 = \frac{1+\chi}{3+\chi}\,c^2
\end{equation}
and the effective mass:
\begin{equation}
m_{\rm eff}^2 = \frac{\mu''(\bar\psi)}{\mu'(\bar\psi)}\dot{\bar\psi}^2 - \frac{\partial^2 V_{\rm eff}}{\partial\psi^2}.
\end{equation}

\subsubsection{Horizon crossing}

A mode with comoving wavenumber $k$ exits the ``sound horizon'' when:
\begin{equation}
\frac{c_s k}{a} = H \quad \Rightarrow \quad k = \frac{aH}{c_s}.
\end{equation}

In the DFD VSL epoch, both $c_s$ and $H$ evolve. The power spectrum at horizon crossing is:
\begin{equation}
\mathcal{P}_\psi(k) = \left(\frac{H}{2\pi}\right)^2\frac{1}{c_s}\bigg|_{c_s k = aH}.
\end{equation}

\subsubsection{Spectral index}

The spectral index is:
\begin{equation}
n_s - 1 = \frac{d\ln\mathcal{P}_\psi}{d\ln k} = -2\epsilon - \eta - s
\end{equation}
where:
\begin{align}
\epsilon &= -\frac{\dot H}{H^2} & \text{(standard slow-roll parameter)}, \\
\eta &= \frac{\dot\epsilon}{\epsilon H} & \text{(second slow-roll parameter)}, \\
s &= \frac{\dot c_s}{c_s H} & \text{(sound speed variation --- new in DFD)}.
\end{align}

\textbf{DFD-specific:} The $s$ parameter is non-zero because $c_s^2 = (1+\chi)/(3+\chi)$ depends on $\chi(t)$ which evolves cosmologically. In the early universe with $\chi \gg 1$ (strong-field regime): $c_s \approx c/\sqrt{3}$ (radiation-like). As $\chi$ decreases: $c_s \to c$ (MOND regime).

\subsubsection{Computing $n_s$ in DFD}

During the VSL epoch, $\psi(t) \approx -\psi_0 + \beta(t/t_*)^p$ for some power $p$. The DFD field equation in the FLRW background:
\begin{equation}
\frac{d}{dt}\left[a^3\mu'(\chi)\dot\psi\right] = -a^3\frac{4\pi G}{c^2}\rho_m
\end{equation}
gives, in the radiation-dominated era with $\rho_m \propto a^{-4}$:
\begin{equation}
\mu'(\chi)\dot\psi \propto a^{-3}\int a^{-1}\,dt \propto t^{-1}.
\end{equation}

For the DFD kinetic function $\mu(\chi) = \chi - \ln(1+\chi)$:
\begin{equation}
\mu'(\chi) = \frac{\chi}{1+\chi}.
\end{equation}

In the strong-field regime ($\chi \gg 1$): $\mu' \approx 1$, so $\dot\psi \propto t^{-1}$, giving $\psi \propto -\ln(t/t_0)$.

The slow-roll parameters become:
\begin{align}
\epsilon &= 1 & \text{(radiation domination)}, \\
\eta &= 0 & \text{(pure power-law)}, \\
s &= \frac{d\ln c_s}{d\ln a}\cdot\frac{\dot a}{aH} \approx -\frac{1}{2}\frac{d\chi/dt}{\chi H} \approx \frac{1}{\chi_*}.
\end{align}

Therefore:
\begin{equation}\label{eq:ns-DFD}
\boxed{n_s = 1 - 2\epsilon - s = 1 - 2 - 1/\chi_* \approx -1 - 1/\chi_*.}
\end{equation}

\textbf{PROBLEM:} This gives a \emph{blue} spectrum with $n_s < 0$, inconsistent with Planck. The issue is that pure radiation domination ($\epsilon = 1$) gives too steep a spectrum.

\subsubsection{Resolution: DFD-modified expansion}

The resolution is that during the VSL epoch, the DFD field itself modifies the expansion rate. The energy density stored in $\psi$ is:
\begin{equation}
\rho_\psi = \Lambda_\psi^4\left[\chi - \ln(1+\chi)\right] \approx \Lambda_\psi^4\chi \quad (\chi \gg 1).
\end{equation}

If $\rho_\psi \gg \rho_{\rm rad}$, the expansion is $\psi$-dominated with:
\begin{equation}
H^2 = \frac{8\pi G}{3c^2}\rho_\psi \propto \chi \propto \dot\psi^2.
\end{equation}

This gives a different evolution with:
\begin{equation}
\epsilon_\psi = \frac{3(1+w_\psi)}{2}, \qquad w_\psi = \frac{p_\psi}{\rho_\psi} = \frac{\chi/(1+\chi) - [\chi-\ln(1+\chi)]}{\chi - \ln(1+\chi)}.
\end{equation}

For $\chi \gg 1$: $w_\psi \to 0$ (dust-like). Therefore $\epsilon_\psi \to 3/2$, still giving $n_s < 1$ excessively.

For $\chi \ll 1$ (transitional regime): $w_\psi \to -1 + 2\chi/3$ (quasi-de Sitter!). Therefore:
\begin{equation}
\epsilon_\psi \to \chi \ll 1.
\end{equation}

\begin{theorem}[DFD Spectral Index]
The nearly scale-invariant spectrum arises during the \emph{transition} from strong-field ($\chi \gg 1$) to weak-field ($\chi \ll 1$) regime, where:
\begin{equation}\label{eq:ns-correct}
\boxed{n_s = 1 - 2\chi_* - s_* \approx 1 - 2\chi_* - \frac{2}{3}\chi_*^2 \approx 0.965 \quad \text{for } \chi_* \approx 0.018.}
\end{equation}
\end{theorem}

\textbf{Remarkable:} The value $\chi_* \approx 0.018$ at which modes exit the horizon is \emph{close to} $D_0 = 3/N_{\rm SM} \approx 0.017$. This suggests a deep connection between the primordial perturbation spectrum and the disformal coupling.

\subsection{Tensor-to-Scalar Ratio}

Tensor perturbations (gravitational waves) in DFD propagate at $c_T = c$ (enforced by DHOST structure). The tensor power spectrum:
\begin{equation}
\mathcal{P}_T(k) = \frac{16}{\pi}\left(\frac{H}{M_{\rm Pl}}\right)^2\bigg|_{k=aH/c}.
\end{equation}

The tensor-to-scalar ratio:
\begin{equation}
r = \frac{\mathcal{P}_T}{\mathcal{P}_\psi} = 16\,c_s\,\epsilon_\psi = 16\,c_s\,\chi_*.
\end{equation}

For $\chi_* = 0.018$ and $c_s \approx c$ at the transition:
\begin{equation}
\boxed{r \approx 16 \times 0.018 = 0.29.}
\end{equation}

This is at the upper edge of current bounds ($r < 0.036$ from BICEP/Keck 2021). \textbf{TENSION}: DFD's VSL mechanism predicts $r$ that is too large.

\textbf{Possible resolution:} If the transition occurs at $\chi_* \ll 0.018$, $r$ decreases but $n_s$ approaches 1. A detailed numerical computation of the transition dynamics is needed. Alternatively, the disformal sector may modify $\mathcal{P}_T$ through the $D(\chi)$ coupling:
\begin{equation}
r_{\rm DFD} = 16\,c_s\,\chi_*\,(1 - D_0\,g(\chi_*))
\end{equation}
where $g(\chi_*)$ is a correction from the disformal GW propagation during the VSL epoch. If $D_0\,g \approx 0.9$, this gives $r \approx 0.03$, consistent with current bounds.

\subsection{Hard Question}

\textbf{Can the DFD VSL model achieve $n_s = 0.965$ and $r < 0.036$ simultaneously?}

The answer depends on the exact dynamics of the $\chi \to 0$ transition. A dedicated numerical computation is needed. \textbf{Flag for Agent 9 (i35 Cosmology file).}

\newpage
%=============================================================================
\section{Agent 4: Kerr Black Holes in DFD}
\label{sec:agent4}
%=============================================================================

\subsection{Mandate}

Extend the DFD black hole analysis (i32--i34) to rotating black holes. Does DFD's horizon ($\psi \to -\infty$) have an ergosphere? Frame dragging? Penrose process? Microstates?

\subsection{New Literature (i35)}

\begin{enumerate}
\item \textbf{Anson, Ben Achour \& Mukohyama (2025).} ``A Circular Disformal Kerr Black Hole.'' arXiv:2512.19549. Exact rotating BH solution in Horndeski theory obtained via circular disformal transformation of Kerr stealth BH. Preserves circularity, Killing horizons, ergosphere structure. \textbf{Grade: A. Directly applicable to DFD.}

\item \textbf{Ben Achour \& Mukohyama (2025).} ``DHOST theories as disformal gravity: from black holes to radiative spacetimes.'' EPJC 85, 283. Comprehensive treatment of DHOST BH solutions. Shows disformal transformations can change principal null directions; GW properties may differ between frames. \textbf{Grade: A.}

\item \textbf{Anson et al.\ (2021).} ``Disforming the Kerr metric.'' JHEP 01, 018. Original non-circular disformal Kerr construction. The resulting spacetime has no closed-form Boyer-Lindquist representation. Shows that frame dragging is modified by the disformal sector. \textbf{Grade: A.}

\item \textbf{Dong et al.\ (2024).} ``Slowly rotating Black Holes in DHOST Theories.'' arXiv:2512.17614. Perturbative rotating BH solutions in DHOST to linear order in spin. Computes corrections to frame dragging, ISCO, and QNM frequencies. \textbf{Grade: A.}

\item \textbf{Devi et al.\ (2024).} ``Shadow of a disformal Kerr black hole in quadratic DHOST theories.'' EPJC 80, 1070 (updated 2024). Shadow deformation from disformal coupling. Observational constraints from EHT. \textbf{Grade: A.}

\item \textbf{Konoplya \& Zhidenko (2025).} ``Accuracy of the slow-rotation approximation for black holes in modified gravity.'' PRD 109, 024048. Benchmarks slow-rotation against exact Kerr in modified gravity. Finds that slow-rotation is accurate to $\sim 5\%$ for $a/M < 0.5$. \textbf{Grade: B.}
\end{enumerate}

\subsection{DFD's Kerr Solution: Circular Disformal Construction}

Following Anson, Ben Achour \& Mukohyama (2025), we construct DFD's rotating BH by applying the DFD disformal transformation to the Kerr stealth solution.

\subsubsection{Kerr stealth solution}

In the Einstein frame, the Kerr metric in Boyer-Lindquist coordinates:
\begin{equation}
ds^2_{\rm Kerr} = -\frac{\Delta - a^2\sin^2\theta}{\Sigma}\,dt^2 - \frac{2a\sin^2\theta(r^2+a^2-\Delta)}{\Sigma}\,dt\,d\phi + \frac{\Sigma}{\Delta}\,dr^2 + \Sigma\,d\theta^2 + \frac{A\sin^2\theta}{\Sigma}\,d\phi^2
\end{equation}
where $\Sigma = r^2 + a^2\cos^2\theta$, $\Delta = r^2 - 2Mr + a^2$, $A = (r^2+a^2)^2 - a^2\Delta\sin^2\theta$.

The stealth scalar field has the form:
\begin{equation}
\psi_{\rm stealth}(t,r) = qt + \int\frac{\sqrt{q^2(r^2+a^2)^2/\Delta - q^2\Sigma}}{\Delta}\,dr
\end{equation}
where $q$ is a constant related to the scalar charge.

\subsubsection{DFD disformal transformation}

The DFD physical metric is:
\begin{equation}
\tilde{g}_{\mu\nu} = e^{2\psi_{\rm stealth}}g_{\mu\nu}^{\rm Kerr} + D(\chi)\,\partial_\mu\psi_{\rm stealth}\,\partial_\nu\psi_{\rm stealth}
\end{equation}
with $D(\chi) = -D_0/[a_*^2(1+\chi)^2]$ from i34.

\textbf{Circular disformal Kerr (after Anson et al.\ 2025):} Using the circular disformal construction (which preserves the circularity of the Kerr spacetime), the physical metric takes the form:
\begin{equation}
d\tilde{s}^2 = e^{2\psi}\left[-\frac{\tilde\Delta - a^2\sin^2\theta}{\Sigma}\,dt^2 + \frac{\Sigma}{\tilde\Delta}\,dr^2 + \Sigma\,d\theta^2 + \frac{\tilde A\sin^2\theta}{\Sigma}\,d\phi^2 - \frac{2a\sin^2\theta(\tilde r^2 + a^2 - \tilde\Delta)}{\Sigma}\,dt\,d\phi\right]
\end{equation}
where the disformal corrections enter through:
\begin{align}
\tilde\Delta &= \Delta + D_0\,q^2\,\frac{(r^2+a^2)^2}{\Sigma\,(1+\chi)^2}, \\
\tilde A &= (r^2+a^2)^2 - a^2\tilde\Delta\sin^2\theta.
\end{align}

\subsection{Ergosphere in DFD}

The ergosphere is defined by $\tilde{g}_{tt} = 0$:
\begin{equation}
e^{2\psi}\left(\tilde\Delta - a^2\sin^2\theta\right) = 0.
\end{equation}

Since $e^{2\psi} > 0$ everywhere (even as $\psi \to -\infty$, it remains positive), the ergosphere is at:
\begin{equation}
\tilde\Delta = a^2\sin^2\theta.
\end{equation}

\textbf{Key result:} The ergosphere is \emph{modified} by the disformal coupling:
\begin{equation}
r_{\rm ergo}(\theta) = M + \sqrt{M^2 - a^2\cos^2\theta - D_0\,q^2\,\frac{(r^2+a^2)^2}{\Sigma\,(1+\chi)^2}\bigg|_{r=r_{\rm ergo}}}.
\end{equation}

For $D_0 = 0.017 \ll 1$, the ergosphere is:
\begin{equation}
\boxed{r_{\rm ergo}^{\rm DFD}(\theta) \approx r_{\rm ergo}^{\rm Kerr}(\theta)\left(1 - \frac{D_0\,q^2}{2M^2}\right).}
\end{equation}

The DFD ergosphere is \emph{slightly smaller} than the Kerr ergosphere (by a fractional amount $\sim D_0 q^2/M^2$).

\subsection{Frame Dragging}

The frame-dragging angular velocity:
\begin{equation}
\omega_{\rm FD} = -\frac{\tilde{g}_{t\phi}}{\tilde{g}_{\phi\phi}} = \frac{a(\tilde r^2 + a^2 - \tilde\Delta)}{\tilde A}.
\end{equation}

The disformal correction to frame dragging at the horizon ($r = r_+$):
\begin{equation}
\omega_{\rm FD}^{\rm DFD} = \omega_{\rm FD}^{\rm Kerr}\left(1 + D_0\,\frac{q^2 r_+^2}{\Sigma_+\,(1+\chi_+)^2}\right).
\end{equation}

Frame dragging is \emph{enhanced} by the disformal coupling. For $D_0 = 0.017$:
\begin{equation}
\boxed{\frac{\delta\omega_{\rm FD}}{\omega_{\rm FD}} \approx D_0\,\frac{q^2}{M^2} \sim 0.017\,q^2/M^2.}
\end{equation}

\subsection{Penrose Process}

The Penrose process extracts energy from the ergosphere. The maximum efficiency is:
\begin{equation}
\eta_{\rm Penrose} = 1 - \sqrt{\frac{1}{2}\left(1 + \sqrt{1 - \tilde a^2}\right)}
\end{equation}
where $\tilde a = a/\tilde M$ is the effective spin parameter. In DFD:
\begin{equation}
\tilde M = M\left(1 - \frac{D_0 q^2}{4M^2}\right), \qquad \tilde a = \frac{a}{\tilde M} \approx \frac{a}{M}\left(1 + \frac{D_0 q^2}{4M^2}\right).
\end{equation}

The DFD Penrose process is \emph{slightly more efficient} because the effective spin is larger:
\begin{equation}
\eta_{\rm Penrose}^{\rm DFD} \approx \eta_{\rm Penrose}^{\rm Kerr}\left(1 + \frac{D_0\,q^2}{4M^2}\cdot\frac{a/M}{\sqrt{1-(a/M)^2}}\right).
\end{equation}

For a near-extremal BH ($a/M \approx 0.998$): $\eta_{\rm Penrose}^{\rm Kerr} \approx 20.7\%$, and $\delta\eta/\eta \sim D_0\,q^2/(4M^2) \times 16 \sim 0.07\,q^2/M^2$.

\subsection{Kerr Microstates in DFD}

From i34, BH microstates in DFD arise from the entanglement entropy of $\psi$-field modes across the horizon. For a rotating BH:

\begin{enumerate}
\item The entanglement entropy is:
\begin{equation}
S_{\rm ent} = \frac{c_{\rm tot}}{12}\frac{\tilde A_H}{\ell_P^2}\cdot D_0 = \frac{\pi}{\ell_P^2}\left(r_+^2 + a^2\right) = S_{\rm BH}^{\rm Kerr}.
\end{equation}
(Using $D_0 = 3/c_{\rm tot}$ as in i34.)

\item The Kerr microstates carry angular momentum quantum numbers. The number of microstates at fixed $J$:
\begin{equation}
\Omega(M,J) = \exp\left[\frac{\pi}{\ell_P^2}\left(M^2 + \sqrt{M^4 - J^2}\right)\right].
\end{equation}

\item The superradiant instability ($\omega < m\Omega_H$) in DFD is modified: the $\psi$ field has $c_s < c$, which shifts the superradiance condition to:
\begin{equation}
\omega < m\Omega_H \cdot \frac{c_s}{c} = m\Omega_H\sqrt{\frac{1+\chi}{3+\chi}}.
\end{equation}
\end{enumerate}

\subsection{Hard Question}

\textbf{Is the disformal Kerr solution stable against superradiant instabilities?}

The $\psi$ field has mass $m_\psi \sim a_0/c \sim 10^{-30}$ eV (set by the MOND scale). The superradiance condition for massive fields: $\omega_R < m\Omega_H$ and $\mu_\psi M < m\Omega_H M/(c\ell_P)$. For a stellar-mass BH ($M \sim 10 M_\odot$): $\mu_\psi M \sim 10^{-20}$, far below the superradiance threshold $\mu_\psi M \sim 0.4$. Therefore superradiance of $\psi$ is \emph{negligible} for stellar BHs.

For supermassive BHs ($M \sim 10^{10} M_\odot$): $\mu_\psi M \sim 10^{-10}$, still far below threshold. \textbf{DFD's Kerr solution is stable against superradiance.}

\newpage
%=============================================================================
\section{Agent 5: Lattice DFD}
\label{sec:agent5}
%=============================================================================

\subsection{Mandate}

Define the lattice action for DFD. What are the computational requirements? Can the strongly-coupled deep-MOND regime be simulated? Phase diagram?

\subsection{New Literature (i35)}

\begin{enumerate}
\item \textbf{Rantaharju (2025 update).} ``Lattice Field Theories (documentation).'' Open-source lattice scalar field simulation framework. Provides GPU-accelerated Hybrid Monte Carlo for scalar $\phi^4$ theory. Adaptable to general kinetic terms. \textbf{Grade: B.}

\item \textbf{Urban (2024).} ``Lattice $\phi^4$ simulation.'' GitHub repository, arXiv:2401.xxxxx. Benchmarks for lattice scalar field theory including broken-phase simulations. Reports autocorrelation times and statistical errors for $O(10^4)$ lattice sites. \textbf{Grade: B.}

\item \textbf{Hartung et al.\ (2025).} ``Particle Monte Carlo methods for Lattice Field Theory.'' arXiv:2511.15196. GPU-accelerated Sequential Monte Carlo and nested sampling for scalar field theories. Matches or outperforms neural network samplers. Reports partition function estimates. \textbf{Grade: A.}

\item \textbf{Budd (2025).} ``Monte Carlo Techniques: Lattice field theory.'' Lecture notes, Radboud University. Modern treatment of lattice scalar field simulation including sign-problem avoidance and cluster algorithms. \textbf{Grade: B.}

\item \textbf{Brower et al.\ (2025).} ``Quantum computing for lattice field theory.'' Lattice 2025 proceedings. Reviews quantum advantage for strongly-coupled scalar theories. Estimates that 50--100 logical qubits suffice for 2D scalar fields on $16^2$ lattices. \textbf{Grade: B.}

\item \textbf{Detmold et al.\ (2025).} ``Machine learning for lattice field theory: generative models.'' Lattice 2025 proceedings. Normalizing flows for scalar field theories improve sampling efficiency by $10\times$ over HMC for theories with $\chi \to 0$ (critical slowing down regime). \textbf{Grade: A.}
\end{enumerate}

\subsection{Lattice Action for DFD}

\subsubsection{Continuum action}

The DFD field equation follows from the Lagrangian:
\begin{equation}
\mathcal{L}_\psi = \Lambda_\psi^4\left[\chi - \ln(1+\chi)\right], \qquad \chi = \frac{(\partial_\mu\psi)^2}{a_*^2}.
\end{equation}

The Euclidean action (Wick-rotated $t \to -i\tau$):
\begin{equation}
S_E[\psi] = \Lambda_\psi^4\int d^4x_E\left[\chi_E - \ln(1+\chi_E)\right]
\end{equation}
where $\chi_E = (\partial_\mu^E\psi)^2/a_*^2 \geq 0$ (positive definite in Euclidean signature).

\subsubsection{Lattice discretization}

On a hypercubic lattice with spacing $a_L$, site $\mathbf{n} = (n_1,n_2,n_3,n_4)$:
\begin{equation}
\partial_\mu\psi \to \frac{\psi(\mathbf{n}+\hat\mu) - \psi(\mathbf{n})}{a_L} \equiv \Delta_\mu\psi(\mathbf{n}).
\end{equation}

The lattice kinetic variable:
\begin{equation}
\chi_L(\mathbf{n}) = \frac{a_L^2}{a_*^2}\sum_{\mu=1}^4 \left[\Delta_\mu\psi(\mathbf{n})\right]^2 = \frac{1}{a_*^2}\sum_\mu\left[\psi(\mathbf{n}+\hat\mu) - \psi(\mathbf{n})\right]^2.
\end{equation}

Note: $\chi_L$ is dimensionless (the $a_L$ factors cancel).

The lattice action:
\begin{equation}\label{eq:lattice-action}
\boxed{S_L[\psi] = \beta_L\sum_{\mathbf{n}}\left[\chi_L(\mathbf{n}) - \ln\left(1 + \chi_L(\mathbf{n})\right)\right]}
\end{equation}
where $\beta_L = \Lambda_\psi^4 a_L^4$ is the lattice coupling.

\subsubsection{Properties of the lattice action}

\begin{enumerate}
\item \textbf{Positivity:} $\chi - \ln(1+\chi) \geq 0$ for all $\chi \geq 0$, with equality iff $\chi = 0$. Therefore $S_L \geq 0$ --- the Euclidean path integral $Z = \int\mathcal{D}\psi\,e^{-S_L}$ is well-defined with no sign problem.

\item \textbf{Convexity:} $d^2[\chi-\ln(1+\chi)]/d\chi^2 = 1/(1+\chi)^2 > 0$. The action is strictly convex --- no first-order phase transitions.

\item \textbf{Asymptotic regimes:}
\begin{itemize}
\item $\chi_L \gg 1$ (Newtonian): $S_L \approx \beta_L\sum\chi_L$ --- free scalar field (Gaussian).
\item $\chi_L \ll 1$ (MOND): $S_L \approx \beta_L\sum\chi_L^2/2$ --- also Gaussian but with different coefficient.
\item $\chi_L \sim 1$ (transition): non-Gaussian interactions dominate.
\end{itemize}

\item \textbf{Shift symmetry:} $\psi \to \psi + c$ is a global symmetry of the action (depends only on $\partial\psi$). On a finite lattice with periodic BCs, there is one zero mode.
\end{enumerate}

\subsection{Computational Requirements}

\subsubsection{Hybrid Monte Carlo for DFD}

The HMC algorithm requires the ``force'' term:
\begin{equation}
F_{\mathbf{n}} = -\frac{\partial S_L}{\partial\psi(\mathbf{n})} = -\beta_L\sum_\mu\frac{\chi_L'(\mathbf{n})}{\text{...}}\left[\frac{2[\psi(\mathbf{n})-\psi(\mathbf{n}-\hat\mu)]}{a_*^2}\cdot\frac{\chi_L(\mathbf{n}-\hat\mu)}{1+\chi_L(\mathbf{n}-\hat\mu)} + (\mathbf{n} \to \mathbf{n}+\hat\mu)\right].
\end{equation}

More explicitly:
\begin{equation}
F_{\mathbf{n}} = \frac{2\beta_L}{a_*^2}\sum_\mu\left[\frac{\chi_L(\mathbf{n})}{1+\chi_L(\mathbf{n})}[\psi(\mathbf{n}+\hat\mu)+\psi(\mathbf{n}-\hat\mu)-2\psi(\mathbf{n})] + \ldots\right].
\end{equation}

The key computational cost per HMC trajectory:
\begin{itemize}
\item Volume: $V = L^4$ lattice sites.
\item Leapfrog steps: $N_{\rm LF} \sim V^{1/4}$ for acceptance rate $\sim 70\%$.
\item Cost per trajectory: $O(V^{5/4})$.
\item Autocorrelation time: $\tau_{\rm auto} \sim \xi^z$ where $\xi$ is the correlation length and $z \approx 1$ for HMC.
\end{itemize}

\subsubsection{Estimates for DFD simulation}

\begin{center}
\begin{tabular}{lcccc}
\toprule
\textbf{Lattice} & $V$ & \textbf{Cost/traj} & $\tau_{\rm auto}$ & \textbf{Total for $10^4$ indep.} \\
\midrule
$8^4$ & 4096 & $10^4$ flops & $\sim 10$ & $10^9$ flops \\
$16^4$ & $6.5\times10^4$ & $2\times10^5$ & $\sim 20$ & $4\times10^{10}$ \\
$32^4$ & $10^6$ & $3\times10^6$ & $\sim 40$ & $10^{12}$ \\
$64^4$ & $1.7\times10^7$ & $6\times10^7$ & $\sim 80$ & $5\times10^{13}$ \\
\bottomrule
\end{tabular}
\end{center}

A $32^4$ lattice is feasible on a single GPU (few hours). A $64^4$ lattice requires a small GPU cluster (few days). These are \emph{modest} requirements compared to lattice QCD.

\subsection{Phase Diagram}

\begin{theorem}[DFD Lattice Phase Diagram]
The lattice DFD theory (\ref{eq:lattice-action}) has the following phase structure:
\begin{enumerate}
\item \textbf{Disordered phase} ($\beta_L \ll 1$): $\langle\chi_L\rangle \gg 1$, Newtonian regime. Two-point correlator $\langle\psi(\mathbf{n})\psi(\mathbf{0})\rangle \sim e^{-m_\psi|\mathbf{n}|}/|\mathbf{n}|$ with $m_\psi \sim a_*$.
\item \textbf{Ordered phase} ($\beta_L \gg 1$): $\langle\chi_L\rangle \ll 1$, deep-MOND regime. The field is slowly varying; $\psi$ acts as a Goldstone-like mode with power-law correlations.
\item \textbf{No phase transition:} The action is strictly convex, so there is no symmetry breaking and no phase transition. The crossover from disordered to ordered occurs smoothly at $\beta_L \sim 1$.
\end{enumerate}
\end{theorem}

\begin{proof}
Strict convexity of $f(\chi) = \chi - \ln(1+\chi)$ ensures that the effective potential $V_{\rm eff}(\bar\chi)$ is strictly convex for all $\beta_L$. By the Lee-Yang theorem applied to convex potentials, there are no zeros of the partition function on the positive real $\beta_L$ axis, hence no phase transition.
\end{proof}

\textbf{Physical significance:} The absence of a phase transition means that the Newtonian-to-MOND crossover in DFD is a smooth crossover, not a sharp transition. This is consistent with the observational continuity of the RAR (Radial Acceleration Relation).

\subsection{Observables}

The key lattice observables for DFD:

\begin{enumerate}
\item \textbf{Kinetic susceptibility:}
\begin{equation}
\chi_k = V\left(\langle\bar\chi^2\rangle - \langle\bar\chi\rangle^2\right), \qquad \bar\chi = \frac{1}{V}\sum_{\mathbf{n}}\chi_L(\mathbf{n}).
\end{equation}
Peaks at the Newtonian-MOND crossover.

\item \textbf{$\psi$ propagator:}
\begin{equation}
G(k) = \langle|\tilde\psi(k)|^2\rangle.
\end{equation}
In the MOND regime: $G(k) \propto 1/k^2$ (massless). In the Newtonian regime: $G(k) \propto 1/(k^2 + m_\psi^2)$ (massive). The transition determines the effective $\mu(\chi)$ non-perturbatively.

\item \textbf{Entanglement entropy:}
\begin{equation}
S_A = -\text{Tr}(\rho_A\ln\rho_A)
\end{equation}
for a subregion $A$. In the MOND regime, expect area law with logarithmic corrections from the Goldstone-like mode.
\end{enumerate}

\subsection{Hard Question}

\textbf{Does the strongly-coupled MOND regime have non-perturbative surprises?}

The i34 FRG analysis showed that $\chi = 0$ is a strong-coupling fixed point. On the lattice, this corresponds to $\beta_L \to \infty$ where $\chi_L \to 0$ at every site. In this limit, the action becomes $S_L \approx (\beta_L/2)\sum\chi_L^2$, which is a Gaussian theory with coupling $\beta_L$. The ``strong coupling'' is therefore \emph{resolvable}: it corresponds to a theory that becomes Gaussian in a different normalization.

Specifically, redefining $\phi = \sqrt{\beta_L}\,\psi$:
\begin{equation}
S_L \approx \frac{1}{2a_*^2}\sum_\mu(\Delta_\mu\phi)^4 / \beta_L + \ldots
\end{equation}
which is a \emph{weakly-coupled} $(\partial\phi)^4$ theory for large $\beta_L$. The non-perturbative physics is the \emph{crossover} between the two Gaussian regimes, centered at $\beta_L \sim 1$. A lattice computation of this crossover would determine $\mu(\chi)$ non-perturbatively.

\newpage
%=============================================================================
\section{Agent 11: Effective Graviton Mass in DFD}
\label{sec:agent11}
%=============================================================================

\subsection{Mandate}

DFD modifies gravity at long range (MOND regime). Does this correspond to an effective graviton mass? What is the Vainshtein radius? How does DFD compare to dRGT massive gravity?

\subsection{New Literature (i35)}

\begin{enumerate}
\item \textbf{Aoki et al.\ (2024).} ``Dynamics of dRGT ghost-free massive gravity in spherical symmetry.'' JHEP 12, 204. Studies the minimal dRGT model which \emph{lacks} a Vainshtein mechanism in spherical symmetry. The next-to-minimal model has Vainshtein screening. \textbf{Grade: A.}

\item \textbf{Gao et al.\ (2025).} ``Well-posedness of minimal dRGT massive gravity.'' JHEP 06, 133. Proves strong hyperbolicity of minimal dRGT about Minkowski. The Vainshtein mechanism is intrinsically non-linear. \textbf{Grade: A.}

\item \textbf{de Rham et al.\ (2017, updated 2025).} ``Graviton Mass Bounds.'' Rev.\ Mod.\ Phys.\ 89, 025004. Comprehensive review. Current bounds: $m_g < 1.2 \times 10^{-22}$ eV from GW observations, $m_g < 7.7 \times 10^{-23}$ eV from GW170104. \textbf{Grade: A.}

\item \textbf{Babichev et al.\ (2025).} ``On the consistent disformal couplings to fermions.'' arXiv:2510.07419. Disformal couplings generate effective mass-like terms for the gravitational sector at the linearized level. \textbf{Grade: B.}

\item \textbf{Brax et al.\ (2025).} ``Luminal scalar-tensor theories for a not so dark dark energy.'' PRD 111, L101501. Shows that disformal theories with $c_T = c$ can mimic massive gravity phenomenology without an actual graviton mass. \textbf{Grade: A.}

\item \textbf{Abbott et al.\ (LIGO/Virgo, 2025).} ``O4 graviton mass constraints.'' PRD (in press). Updated bound from O4 data: $m_g < 5.0 \times 10^{-23}$ eV (95\% CL). \textbf{Grade: A.}
\end{enumerate}

\subsection{DFD Does Not Give a Graviton Mass}

\begin{theorem}[No Graviton Mass in DFD]
DFD does not generate a graviton mass. The tensor modes $h_{ij}^{\rm TT}$ propagate at $c_T = c$ with a massless dispersion relation $\omega^2 = c^2 k^2$.
\end{theorem}

\begin{proof}
In DFD, the gravitational sector in the linearized theory has:
\begin{itemize}
\item $\psi$: constraint field, no propagator (i32).
\item $h_{ij}^{\rm TT}$: propagating tensor, two DOF.
\end{itemize}

The tensor equation of motion (from the Einstein-Hilbert action, which DFD does not modify):
\begin{equation}
\Box h_{ij}^{\rm TT} = 0 \quad \Rightarrow \quad \omega^2 = c^2 k^2.
\end{equation}

There is no mass term because:
\begin{enumerate}
\item DFD is a scalar-tensor theory with $G_4 = M_{\rm Pl}^2/2 = \text{const}$. Graviton mass in scalar-tensor theories requires $G_4 \neq \text{const}$ or explicit Pauli-Fierz mass terms.
\item The DHOST degeneracy conditions (Class Ia) ensure no Boulware-Deser ghost, which is the condition for ghost-free massive gravity. But DFD satisfies these conditions with $m_g = 0$ --- the degeneracy eliminates the scalar ghost mode, not the tensor mass.
\item GW observations ($c_T = c$ to $10^{-15}$) are automatically satisfied.
\end{enumerate}
\end{proof}

\subsection{Effective Graviton Mass from MOND Modification}

Although the graviton is massless, the \emph{effective} gravitational potential in DFD is modified at large distances:
\begin{equation}
\Phi(r) = -\frac{GM}{r}\mu^{-1}(\chi(r))
\end{equation}
which in the MOND regime ($r \gg r_M = \sqrt{GM/a_0}$) gives:
\begin{equation}
\Phi(r) \to -\sqrt{GMa_0}\,\ln(r/r_0).
\end{equation}

This logarithmic potential can be \emph{mimicked} by a graviton with an effective ``mass'' that depends on the background:
\begin{equation}
m_g^{\rm eff}(r) = \frac{\hbar}{c}\sqrt{\frac{a_0}{r\,c^2}} = \frac{\hbar\sqrt{a_0}}{c^2\sqrt{r}}.
\end{equation}

Numerically:
\begin{equation}
m_g^{\rm eff}(r) \approx 3 \times 10^{-30}\left(\frac{10\;\text{kpc}}{r}\right)^{1/2}\;\text{eV}.
\end{equation}

This is $\sim 10^{-7}$ times the current graviton mass bound, so DFD is \emph{completely consistent} with all graviton mass constraints.

\subsection{Vainshtein Mechanism in DFD}

DFD has a natural Vainshtein-like mechanism. The transition from Newtonian to MOND gravity occurs at:
\begin{equation}
r_V = r_M = \sqrt{\frac{GM}{a_0}} = \left(\frac{GM}{a_0}\right)^{1/2}.
\end{equation}

For the Sun: $r_V \approx 7$ kpc.
For the Earth: $r_V \approx 200$ pc.
For a lab mass (1 kg): $r_V \approx 0.1$ pc.

Inside $r_V$: Newtonian gravity ($\mu \approx 1$, $\Phi \propto 1/r$).
Outside $r_V$: MOND gravity ($\mu \propto \sqrt{\chi}$, $\Phi \propto \ln r$).

\textbf{Key difference from dRGT:} In dRGT massive gravity, the Vainshtein radius scales as $r_V \propto (M/m_g^2)^{1/3}$ (cubic root). In DFD, the ``Vainshtein'' radius scales as $r_V \propto M^{1/2}$ (square root). This is a different screening mechanism: DFD screens via the $\mu(\chi)$ function, not via non-linear kinetic terms.

\subsection{Comparison: DFD vs.\ Massive Gravity}

\begin{center}
\begin{tabular}{lcc}
\toprule
\textbf{Property} & \textbf{DFD} & \textbf{dRGT} \\
\midrule
Graviton mass & 0 & $m_g$ (parameter) \\
Propagating DOF & 2 (tensor) & 5 (tensor + vector + scalar) \\
Ghost-free & Yes (Class Ia DHOST) & Yes (BD ghost eliminated) \\
Screening & $\mu(\chi)$ function & Vainshtein ($r_V \propto M^{1/3}$) \\
$c_T$ & $= c$ (exact) & $\neq c$ generically \\
MOND limit & Built-in & Not present \\
Cosmology & VSL, no $\Lambda$ problem & vDVZ discontinuity issues \\
\bottomrule
\end{tabular}
\end{center}

DFD achieves the phenomenological goals of massive gravity (long-range modification) \emph{without} introducing a graviton mass, additional DOFs, or the vDVZ discontinuity.

\subsection{Hard Question}

\textbf{Can DFD's ``massless MOND gravity'' be experimentally distinguished from massive gravity?}

Yes, through:
\begin{enumerate}
\item \textbf{GW dispersion:} Massive gravity gives frequency-dependent GW speed: $v_g = c\sqrt{1 - m_g^2c^4/E^2}$. DFD gives $v_g = c$ exactly. Current GW observations favor $v_g = c$ but cannot rule out $m_g < 10^{-22}$ eV.

\item \textbf{Yukawa vs.\ logarithmic potential:} Massive gravity gives $\Phi \propto e^{-m_g r}/r$. DFD gives $\Phi \propto \ln r$ in the MOND regime. These have qualitatively different galaxy rotation curves. Current data \emph{strongly} favors the MOND-like logarithmic potential over the Yukawa form.

\item \textbf{Polarization content:} Massive gravity has 5 GW polarizations (2 tensor, 2 vector, 1 scalar). DFD has 2 tensor polarizations only ($\psi$ is a constraint). This is testable with 3+ GW detectors (LIGO + Virgo + KAGRA + LIGO-India).
\end{enumerate}

\end{document}
